\subsection{Solutions des ateliers}
\begin{enumerate}
\item Les sorties sont $(f\circ g)(x)=x^{2} + 4 x + 3$ et $(g\circ f)(x)=x^{2} + 1$. Leur différence est $4 x + 2$, nulle uniquement pour $x=- \frac{1}{2}$.
\item Sur $[-2,2]$, la distance à 1 varie de 0 à 3 ; ainsi $h([-2,2])=[2,11]$. La condition $3\leq h(x)\leq6$ équivaut à $1\leq|x-1|\leq2$ ; l'image réciproque est $[-1,0]\cup[2,3]$. Les deux singletons A et C sont disjoints, mais ont la même image $\{6\}$ : l'image de leur intersection est vide, tandis que l'intersection de leurs images ne l'est pas.
\item La facture finale vaut
\[q(x)=\begin{cases}x+2&0\leq x\leq2,\\2x&2<x\leq5/2,\\x+5/2&5/2<x\leq4.\end{cases}\]
Le changement de tarif se fait à deux heures ; la réduction débute lorsque $2x=5$. Les raccords valent respectivement 4 et 5 des deux côtés. Les deux premières branches restent au plus égales à 5 ; une facture de 6 relève donc de la dernière branche et donne $x=7/2$. La substitution dans les deux traitements redonne 6.
\end{enumerate}
